We view our
theoretical framework for sharpening as a starting point toward a foundational
understanding of self-improvement that can guide the design and evaluation of algorithms. To this end, \abedit{we raise several} directions
for future research.
\begin{itemize}
\item \emph{Representation learning.} A conceptually appealing feature
  of our framework is that it is agnostic to the
  structure of the model under consideration, but an important
  direction for future work is to study the dynamics of
  self-improvement for specific models/architectures and
  understand the representations that these models learn under
  self-training.
\item \emph{Richer forms of self-reward.} Our theoretical results study the
  dynamics of self-training in a stylized framework where the model
  uses its own log-probabilities as a self-reward. Empirical research on
  self-improvement leverages more sophisticated \abedit{approaches} (e.g., specific prompting
  techniques)
  \citep{huang2022large,wang2022self,bai2022constitutional,pang2023language,yuan2024self}
  and it is important to understand when and how these forms of
  self-improvement are beneficial.
\end{itemize}



